\section{Next Smaller Element}
\label{sec:sq-next-smaller}

\problemstatement{Given an array $\textit{arr}$, find, for every element,
the next element to its right that is strictly \textbf{smaller}, or $-1$
if none exists.}

\begin{patternbox}
The mirror image of Section~\ref{sec:sq-next-greater-1}: flip every
comparison, and the monotonic stack should now be kept strictly
\textbf{increasing} from bottom to top instead of decreasing.
\end{patternbox}

\subsection*{Understanding the problem carefully}

By the same reasoning as next greater element, scanning right to left:
before processing $\textit{arr}[i]$, any stack element $\ge
\textit{arr}[i]$ can never be the ``next smaller'' answer for anything
still to the left, since $\textit{arr}[i]$ is closer and at least as
small. Popping these and keeping whatever survives gives the nearest
qualifying smaller value.

\begin{ideabox}[Monotonic Increasing Stack, Scanned Right to Left]
Scan right to left. Before processing $\textit{arr}[i]$, pop every stack
element $\ge \textit{arr}[i]$. The survivor on top (if any) is $i$'s next
smaller element. Push $\textit{arr}[i]$.
\end{ideabox}

\subsection*{Algorithm}

\begin{enumerate}
  \item Initialize an empty stack and a result array.
  \item For $i$ from $n-1$ down to $0$:
  \begin{enumerate}
    \item While the stack is not empty and its top $\ge
          \textit{arr}[i]$: pop.
    \item $\textit{result}[i] \gets$ the stack's top if non-empty, else
          $-1$.
    \item Push $\textit{arr}[i]$.
  \end{enumerate}
  \item Return \textit{result}.
\end{enumerate}

\subsection*{Dry run}

\dryrun{Take $\textit{arr} = [4,8,5,2,25]$.}

\begin{itemize}
  \item $i=4$ ($25$): stack empty. $\textit{result}[4]=-1$. Push $25$.
  \item $i=3$ ($2$): top $25\ge2$, pop. Stack empty.
        $\textit{result}[3]=-1$. Push $2$.
  \item $i=2$ ($5$): top $2<5$, no pop. $\textit{result}[2]=2$. Push
        $5$. Stack $=[2,5]$.
  \item $i=1$ ($8$): top $5<8$, no pop. $\textit{result}[1]=5$. Push
        $8$. Stack $=[2,5,8]$.
  \item $i=0$ ($4$): top $8\ge4$, pop. New top $5\ge4$, pop. New top
        $2<4$, no pop. $\textit{result}[0]=2$. Push $4$.
\end{itemize}

Result: $[2,5,2,-1,-1]$.

\subsection*{C++ implementation}

\begin{lstlisting}
#include <bits/stdc++.h>
using namespace std;

vector<int> nextSmallerElement(vector<int>& arr) {
    int n = arr.size();
    vector<int> result(n, -1);
    stack<int> st;

    for (int i = n - 1; i >= 0; i--) {
        while (!st.empty() && st.top() >= arr[i]) {
            st.pop();
        }
        result[i] = st.empty() ? -1 : st.top();
        st.push(arr[i]);
    }
    return result;
}

int main() {
    vector<int> arr = {4, 8, 5, 2, 25};
    for (int x : nextSmallerElement(arr)) cout << x << " ";
    cout << endl;
    return 0;
}
\end{lstlisting}

\begin{complexitybox}
\textbf{Time Complexity.} $O(n)$, by the same amortized argument as
Section~\ref{sec:sq-next-greater-1}.
\textbf{Space Complexity.} $O(n)$.
\end{complexitybox}

\begin{takeawaybox}
``Next/previous greater/smaller'' is really one template with four
dials: which direction to scan (left-to-right for \emph{previous},
right-to-left for \emph{next}) and which comparison the stack enforces
(decreasing for \emph{greater}, increasing for \emph{smaller}). Once you
have derived one of the four, the rest follow by flipping the relevant
dial.
\end{takeawaybox}
